\documentclass[11pt]{article}
\usepackage[margin=1in]{geometry}
\usepackage{amsmath, amssymb, amsthm}
\usepackage{graphicx}
\usepackage{natbib}
\usepackage{hyperref}

\title{Bayesian Learning of Switching Linear Dynamical Systems\\
  \large Statistical Foundations of AI/ML Project}
\author{Krish}
\date{\today}

\begin{document}
\maketitle

\begin{abstract}
Many real-world time series exhibit regime changes, where the underlying
generative dynamics switch among a set of simpler modes. Classical models
such as hidden Markov models (HMMs) and linear dynamical systems (LDSs)
capture either discrete or continuous structure, but not both. In this
project I study switching linear dynamical systems (SLDS) and closely
related switching vector autoregressive (AR) processes, with an emphasis
on Bayesian modeling and inference. Building on the work of
\citet{fox2008} and \citet{linderman2017}, I implemented and empirically
compared three families of models: finite switching AR models,
nonparametric sticky HDP-AR-HMMs, and recurrent AR-HMMs with
state-dependent transitions. Experiments on synthetic and real time
series highlight when nonparametric priors help discover the number of
regimes and when recurrent, state-dependent transitions are necessary to
capture non-geometric dwell-time structure.
\end{abstract}

\section{Introduction}

Complex temporal phenomena in science and engineering often display
structured, recurring modes of behaviour. Examples range from financial
markets alternating between low- and high-volatility regimes to animals
switching between stereotyped motion patterns. A key challenge for
statistical modeling is to represent both the continuous evolution
within modes and the discrete switches between modes.

Linear dynamical systems (LDSs) provide a tractable, Gaussian model for
continuous-valued time series, while hidden Markov models (HMMs) capture
discrete regime switching with geometric state durations. A natural
hybrid is the switching linear dynamical system (SLDS), in which a
latent Markov chain selects among a library of linear-Gaussian dynamical
systems. Closely related are switching vector autoregressive models,
which can be viewed as SLDSs where the latent continuous state is fully
observed. These models have been widely used in control, econometrics
and computer vision.

However, several limitations motivate a more Bayesian perspective.
First, standard SLDS formulations fix the number of modes \emph{a
priori}, and miss the opportunity to let the data inform model
complexity. Second, Markovian switching implies geometric state
durations, which can be unrealistic in applications with long, structured
episodes. Third, classical maximum-likelihood training does not quantify
uncertainty over parameters and latent segmentations.

In this project I explore Bayesian solutions to these issues, focusing on
two influential lines of work:
\begin{itemize}
  \item \textbf{Nonparametric Bayesian SLDS/AR-HMM} of \citet{fox2008},
    which use hierarchical Dirichlet processes (HDPs) and a ``sticky''
    prior to infer an unknown number of persistent dynamical modes.
  \item \textbf{Recurrent SLDS (rSLDS)} of \citet{linderman2017}, which
    allow transition probabilities to depend on the continuous latent
    state via a logistic stick-breaking regression, enabling
    state-dependent, non-geometric dwell times.
\end{itemize}

My contributions are purely algorithmic and empirical: I implemented
finite switching AR models, a sticky HDP-AR-HMM with weak-limit
truncation, and an approximate recurrent AR-HMM using multinomial
logistic regression, and compared them on synthetic and real data.

\section{Background}

\subsection{HMMs, LDSs and SLDSs}

An HMM consists of a discrete Markov chain $(z_t)_{t=1}^T$ with
transition matrix $\Pi = (\pi_{ij})$ and initial distribution $\pi_0$,
and observations $y_t$ drawn conditionally independently given $z_t$,
e.g.\ $y_t \mid z_t = k \sim \mathcal{N}(\mu_k, \Sigma_k)$. The Markov
assumption yields geometric state durations.

An LDS (state-space model) posits a continuous latent state $x_t \in
\mathbb{R}^n$ evolving linearly with Gaussian noise,
\[
  x_t = A x_{t-1} + e_t, \quad e_t \sim \mathcal{N}(0, \Sigma),
  \qquad
  y_t = C x_t + w_t, \quad w_t \sim \mathcal{N}(0, R).
\]
Inference is tractable via Kalman filtering and smoothing.

A switching LDS combines these ideas: a discrete Markov chain $z_t$
selects among a set of linear dynamics and/or observation models
indexed by $k = 1,\dots,K$,
\[
  x_t = A^{(z_t)} x_{t-1} + e_t^{(z_t)},
  \qquad
  y_t = C^{(z_t)} x_t + w_t^{(z_t)}.
\]
In this project I mainly work with the equivalent switching VAR(1)
representation where the state is taken to be the observation itself,
leading to an autoregressive HMM (AR-HMM)
\citep{fox2008}. For each mode $k$,
\[
  y_t = A^{(k)} y_{t-1} + e_t^{(k)}, \quad e_t^{(k)} \sim
  \mathcal{N}(0, \Sigma^{(k)}).
\]

\subsection{Dirichlet processes and the sticky HDP-HMM}

The Dirichlet process (DP) $\mathrm{DP}(\gamma, H)$ is a distribution
over discrete probability measures $G$ on a parameter space $\Theta$,
often constructed via stick-breaking:
\[
  \beta_k' \sim \mathrm{Beta}(1,\gamma),\quad
  \beta_k = \beta_k' \prod_{\ell<k} (1 - \beta_\ell'),\quad
  \theta_k \sim H, \quad
  G = \sum_{k=1}^\infty \beta_k \delta_{\theta_k}.
\]
The hierarchical Dirichlet process (HDP) extends this to multiple
groups, yielding the HDP-HMM when group-specific measures index
rows of an HMM transition matrix \citep{teh2006, beal2002}.

While the HDP-HMM allows an unbounded number of HMM states, it can
over-fragment time series data by rapidly switching among redundant
states. \citet{fox2008} proposed the sticky HDP-HMM, which adds a
``stickiness'' parameter $\kappa > 0$ that increases the prior
probability of self-transitions. In the weak-limit finite
approximation, with truncation level $L$, this leads to
\[
  \beta \sim \mathrm{Dir}\Big(\tfrac{\gamma}{L},\dots,\tfrac{\gamma}{L}\Big),
  \quad
  \pi_j \sim \mathrm{Dir}(\alpha \beta + \kappa \delta_j),
\]
where $\delta_j$ is a point mass on the $j$-th component.

\subsection{Recurrent SLDS and state-dependent transitions}

Standard SLDS and HDP-HMM models assume that transition probabilities
depend only on the previous discrete state $z_{t-1}$. As a result,
state durations are geometric. In many dynamical systems, however, the
probability of switching regimes depends strongly on the continuous
state $x_t$ or observed data $y_t$. For example, a dynamical system may
switch when it reaches a particular region of state space.

\citet{linderman2017} introduced a recurrent SLDS (rSLDS) in which the
transition probabilities are governed by a logistic stick-breaking
regression on the continuous state, enabling rich, state-dependent
switching while maintaining tractable Bayesian inference via P{\'o}lya
gamma augmentation. In this project I approximate this idea at the
level of an AR-HMM by using multinomial logistic regression to model
$p(z_{t+1}\mid z_t, y_t)$.

\section{Models and Inference}

Due to space, I summarize the implemented models and defer derivations
to the source code and lecture notes.

\subsection{Finite switching AR(1) model}

The finite model assumes a fixed number $K$ of modes with Markov
transitions. I implemented both an EM-style maximum-likelihood version
and a Bayesian version with simple conjugate updates for $A^{(k)}$ and
$\Sigma^{(k)}$. Inference over the latent sequence $z_{1:T}$ uses the
forward--backward algorithm and forward-filter backward-sample.

\subsection{Sticky HDP-AR-HMM}

The sticky HDP-AR-HMM follows \citet{fox2008} and augments each mode
with a VAR(1) parameter pair $(A^{(k)}, \Sigma^{(k)})$. I use the
weak-limit truncation at level $L$ and a blocked Gibbs sampler that
alternates between sampling the mode sequence $z_{1:T}$, the
transition distributions $\pi_j$, the global weights $\beta$, and the
AR parameters for each mode. This yields a nonparametric prior over the
number of occupied modes.

\subsection{Recurrent AR-HMM}

To approximate recurrent SLDS behavior without implementing full
P{\'o}lya--gamma augmentation, I model state-dependent transitions via
a multinomial logistic regression classifier mapping features
$(z_t, y_t)$ to $z_{t+1}$. I then alternate between approximate
forward--backward inference using the current transition model,
updating AR dynamics, and refitting the logistic regression on sampled
state sequences. While not fully Bayesian, this captures the key
qualitative behavior of state-dependent, non-geometric switching.

\section{Experiments}

\subsection{Synthetic Markov-switching AR(1)}

I first generated synthetic data from a $K=3$ Markov-switching AR(1)
model with stable dynamics and high self-transition probability
($0.9$). I then fit:
\begin{enumerate}
  \item an EM-style finite AR-HMM with $K=3$,
  \item a Bayesian switching AR(1) model with Gibbs sampling, and
  \item a sticky HDP-AR-HMM with truncation $L=6$.
\end{enumerate}

Performance was evaluated using Hamming distance between the true and
inferred state sequences (optimized over label permutations) and
summary statistics of state run-lengths. Both finite models recovered
the correct segmentation well. The HDP-AR-HMM typically used close to
three effective modes, occasionally splitting a true mode into two when
the data supported subtle dynamical differences, illustrating the
flexibility of the nonparametric prior.

\subsection{Synthetic state-dependent switching on an oval track}

Next, I considered a simplified analogue of the NASCAR-style example
in \citet{linderman2017}, generating a 2D trajectory along an oval
``track'' with four position-dependent modes. The true next mode is
almost deterministic given the current position, producing highly
structured, non-geometric dwell times.

I compared a Markov AR-HMM baseline to the recurrent AR-HMM. The
Markov model captured local dynamics within modes but produced
geometric-like dwell times and sometimes switched in the wrong regions.
In contrast, the recurrent model learned position-dependent transition
boundaries and generated trajectories with much more realistic periodic
structure and dwell-time statistics similar to the ground truth.

\subsection{Real-data example: stock index volatility regimes}

Finally, I applied the finite AR-HMM and sticky HDP-AR-HMM to daily log
returns of a stock index. Both models revealed interpretable
low-, medium- and high-volatility regimes. The sticky HDP-AR-HMM
showed more persistent regimes and occasionally split volatile periods
into sub-regimes with different autocorrelation structure. Run-length
statistics suggested that the nonparametric model better captured the
heavy-tailed nature of volatility episode durations, consistent with
the motivation for sticky HDP priors.

\section{Discussion and Conclusion}

This project explored Bayesian modeling of regime-switching time series,
highlighting three ideas: (i) switching AR(1) models as a simple and
practical alternative to full SLDSs, (ii) hierarchical Dirichlet
process priors and the sticky HDP-HMM for inferring the number of
persistent dynamical modes, and (iii) recurrent, state-dependent
transition models inspired by rSLDS. Synthetic experiments confirmed
known advantages of HDP-based models in controlling over-segmentation
and of recurrent models in capturing non-geometric dwell times. The
real-data results, while preliminary, demonstrated that these models
can discover meaningful volatility regimes in financial time series.

From a statistical foundations perspective, the project connected
Dirichlet processes, conjugate priors for linear Gaussian models,
blocked Gibbs sampling, and logistic regression-based transition
models. Future work could replace the approximate recurrent AR-HMM
with a fully Bayesian rSLDS using P{\'o}lya--gamma augmentation as in
\citet{linderman2017}, and apply the models to higher-dimensional
neural or motion capture datasets.

\bibliographystyle{plainnat}
\begin{thebibliography}{9}

\bibitem[Beal et~al.(2002)Beal, Ghahramani, and Rasmussen]{beal2002}
M.~J. Beal, Z.~Ghahramani, and C.~E. Rasmussen.
\newblock The infinite hidden {M}arkov model.
\newblock In \emph{Advances in Neural Information Processing Systems}, 2002.

\bibitem[Fox et~al.(2008)Fox, Sudderth, Jordan, and Willsky]{fox2008}
E.~B. Fox, E.~B. Sudderth, M.~I. Jordan, and A.~S. Willsky.
\newblock Nonparametric {B}ayesian learning of switching linear dynamical
  systems.
\newblock In \emph{Advances in Neural Information Processing Systems}, 2008.

\bibitem[Linderman et~al.(2017)Linderman, Johnson, Miller, Adams, Blei, and
  Paninski]{linderman2017}
S.~W. Linderman, M.~J. Johnson, A.~C. Miller, R.~P. Adams, D.~M. Blei, and
  L.~Paninski.
\newblock {B}ayesian learning and inference in recurrent switching linear
  dynamical systems.
\newblock In \emph{Proceedings of AISTATS}, 2017.

\bibitem[Teh et~al.(2006)Teh, Jordan, Beal, and Blei]{teh2006}
Y.~W. Teh, M.~I. Jordan, M.~J. Beal, and D.~M. Blei.
\newblock Hierarchical {D}irichlet processes.
\newblock \emph{Journal of the American Statistical Association},
  101\penalty0 (476):\penalty0 1566--1581, 2006.

\end{thebibliography}

\end{document}

